% =========================================================
%  sections/02-cpl.tex
%  Bagian 2 — Capaian Pembelajaran Lulusan (CPL)
% =========================================================

\section{Capaian Pembelajaran Lulusan (CPL)}

Capaian Pembelajaran Lulusan (CPL) adalah kompetensi lulusan yang dibebankan
pada mata kuliah ini. Rumusan CPL mengacu pada profil lulusan Program Studi
\NamaProdi\ dan dinyatakan dalam bentuk kompetensi lulusan yang mencakup aspek
pengetahuan dan keterampilan.

\begin{table}[htbp]
\centering
\caption{CPL yang dibebankan pada mata kuliah \NamaMK}
\label{tab:cpl}
\begin{tabularx}{\textwidth}{@{}C{2.0cm}Y C{3.0cm} C{1.8cm}@{}}
\toprule
\headrow \textbf{Kode} & \textbf{Rumusan CPL} & \textbf{Aspek Kompetensi} & \textbf{Ranah Bloom} \\
\midrule
\CPL{03} &
Kemampuan merancang dan menganalisis algoritma untuk menyelesaikan
permasalahan organisasi secara optimal, serta memilih dan menerapkannya pada
bahasa pemrograman tertentu. &
Pengetahuan (P) dan Keterampilan Umum (KU) &
C4--C5 \\
\addlinespace
\CPL{04} &
Kemampuan mendesain, mengimplementasi, dan mengevaluasi solusi berbasis
\istilah{computing} yang memenuhi kebutuhan-kebutuhan \istilah{computing} pada
sebuah disiplin program. &
Keterampilan Khusus (KK) &
C5--C6 \\
\bottomrule
\end{tabularx}
\end{table}

\subsection*{Penjabaran CPL pada Mata Kuliah}

\begin{itemize}
  \item \textbf{\CPL{03} --- Algoritma dan bahasa pemrograman.}
        Pada mata kuliah ini CPL03 diwujudkan melalui perancangan dan analisis
        algoritma permainan: siklus \istilah{game loop}, manajemen waktu,
        gerak berbasis vektor, fisika sederhana, deteksi tabrakan, serta
        algoritma AI (perilaku \istilah{chase}, FSM, dan A*). Seluruh algoritma
        tersebut diimplementasikan langsung pada bahasa pemrograman Python
        dengan \istilah{library} Pygame, sehingga mahasiswa memahami mekanisme
        internal permainan tanpa bergantung pada fitur \istilah{black-box}
        \istilah{game engine}.

  \item \textbf{\CPL{04} --- Desain, implementasi, dan evaluasi solusi
        \istilah{computing}.}
        Pada mata kuliah ini CPL04 diwujudkan melalui desain dan implementasi
        solusi permainan pada \istilah{game engine} Unity (penyusunan
        \istilah{project} dan \istilah{scene}, pembangunan level, konfigurasi
        fisika dan \istilah{collider}, integrasi aset visual dan audio), serta
        evaluasi solusi tersebut melalui \istilah{playtesting}, pencatatan
        \istilah{bug}, dan perbaikan berulang terhadap prototipe.
\end{itemize}

Kedua CPL tersebut bersifat saling melengkapi: CPL03 membangun kemampuan
algoritmik yang bersifat mendasar, sedangkan CPL04 membangun kemampuan
mendesain dan mengevaluasi produk \istilah{computing} yang utuh. Bobot
kontribusi masing-masing CPL terhadap nilai akhir mata kuliah diuraikan pada
Bagian~\ref{sec:penilaian}.
